% The control-flow graph of the reaching-definitions example (Lesson 10).
% A loop over x[0..n-1] that sums absolute values: B1 initialises i and s,
% B2 loads t, B3 adds a positive t, B4 negates t and adds it, B5
% increments i and loops back to B2 or leaves to EXIT.  Definitions are
% numbered d1..d7.
% Original (the book's own CFG; not the outline's m10l4_gnk2 example).
% Usage: % The control-flow graph of the reaching-definitions example (Lesson 10).
% A loop over x[0..n-1] that sums absolute values: B1 initialises i and s,
% B2 loads t, B3 adds a positive t, B4 negates t and adds it, B5
% increments i and loops back to B2 or leaves to EXIT.  Definitions are
% numbered d1..d7.
% Original (the book's own CFG; not the outline's m10l4_gnk2 example).
% Usage: % The control-flow graph of the reaching-definitions example (Lesson 10).
% A loop over x[0..n-1] that sums absolute values: B1 initialises i and s,
% B2 loads t, B3 adds a positive t, B4 negates t and adds it, B5
% increments i and loops back to B2 or leaves to EXIT.  Definitions are
% numbered d1..d7.
% Original (the book's own CFG; not the outline's m10l4_gnk2 example).
% Usage: % The control-flow graph of the reaching-definitions example (Lesson 10).
% A loop over x[0..n-1] that sums absolute values: B1 initialises i and s,
% B2 loads t, B3 adds a positive t, B4 negates t and adds it, B5
% increments i and loops back to B2 or leaves to EXIT.  Definitions are
% numbered d1..d7.
% Original (the book's own CFG; not the outline's m10l4_gnk2 example).
% Usage: \input{figures/l10-rd-cfg}
% Width: body (120 mm); the drawing is about 92 mm wide.
\begin{tikzpicture}[x=1mm, y=1mm, font=\scriptsize\sffamily,
    >={Stealth[length=1.6mm]},
    bb/.style={draw, fill=ln-box, align=left, inner sep=2pt, font=\scriptsize},
    nm/.style={font=\scriptsize\sffamily\bfseries, inner sep=1pt}]
  \node[text=ln-muted] (en) at (45,66) {ENTRY};
  \node[bb] (b1) at (45,57) {$d_1$: \texttt{i = 0}\\$d_2$: \texttt{s = 0}};
  \node[bb, fill=ln-accent!15] (b2) at (45,42) {$d_3$: \texttt{t = x[i]}};
  \node[bb] (b3) at (22,27) {$d_4$: \texttt{s = s + t}};
  \node[bb] (b4) at (68,27) {$d_5$: \texttt{t = -t}\\$d_6$: \texttt{s = s + t}};
  \node[bb] (b5) at (45,12) {$d_7$: \texttt{i = i + 1}};
  \node[text=ln-muted] (ex) at (45,1) {EXIT};
  \node[nm, anchor=east] at (b1.west) {B1\,};
  \node[nm, anchor=east] at (b2.west) {B2\,};
  \node[nm, anchor=east] at (b3.west) {B3\,};
  \node[nm, anchor=west] at (b4.east) {\,B4};
  \node[nm, anchor=east] at (b5.west) {B5\,};
  \draw[->] (en) -- (b1);
  \draw[->] (b1) -- (b2);
  \draw[->] (b2) -- node[above left, inner sep=1pt] {\texttt{t > 0}} (b3);
  \draw[->] (b2) -- node[above right, inner sep=1pt] {\texttt{t <= 0}} (b4);
  \draw[->] (b3) -- (b5);
  \draw[->] (b4) -- (b5);
  \draw[->] (b5) -- node[right, inner sep=1.5pt] {\texttt{i >= n}} (ex);
  \draw[->, ln-caution, thick] (b5.east) -- ++(40,0) |- node[pos=0.25, right, inner sep=1.5pt]
    {\texttt{i < n}} (b2.east);
\end{tikzpicture}

% Width: body (120 mm); the drawing is about 92 mm wide.
\begin{tikzpicture}[x=1mm, y=1mm, font=\scriptsize\sffamily,
    >={Stealth[length=1.6mm]},
    bb/.style={draw, fill=ln-box, align=left, inner sep=2pt, font=\scriptsize},
    nm/.style={font=\scriptsize\sffamily\bfseries, inner sep=1pt}]
  \node[text=ln-muted] (en) at (45,66) {ENTRY};
  \node[bb] (b1) at (45,57) {$d_1$: \texttt{i = 0}\\$d_2$: \texttt{s = 0}};
  \node[bb, fill=ln-accent!15] (b2) at (45,42) {$d_3$: \texttt{t = x[i]}};
  \node[bb] (b3) at (22,27) {$d_4$: \texttt{s = s + t}};
  \node[bb] (b4) at (68,27) {$d_5$: \texttt{t = -t}\\$d_6$: \texttt{s = s + t}};
  \node[bb] (b5) at (45,12) {$d_7$: \texttt{i = i + 1}};
  \node[text=ln-muted] (ex) at (45,1) {EXIT};
  \node[nm, anchor=east] at (b1.west) {B1\,};
  \node[nm, anchor=east] at (b2.west) {B2\,};
  \node[nm, anchor=east] at (b3.west) {B3\,};
  \node[nm, anchor=west] at (b4.east) {\,B4};
  \node[nm, anchor=east] at (b5.west) {B5\,};
  \draw[->] (en) -- (b1);
  \draw[->] (b1) -- (b2);
  \draw[->] (b2) -- node[above left, inner sep=1pt] {\texttt{t > 0}} (b3);
  \draw[->] (b2) -- node[above right, inner sep=1pt] {\texttt{t <= 0}} (b4);
  \draw[->] (b3) -- (b5);
  \draw[->] (b4) -- (b5);
  \draw[->] (b5) -- node[right, inner sep=1.5pt] {\texttt{i >= n}} (ex);
  \draw[->, ln-caution, thick] (b5.east) -- ++(40,0) |- node[pos=0.25, right, inner sep=1.5pt]
    {\texttt{i < n}} (b2.east);
\end{tikzpicture}

% Width: body (120 mm); the drawing is about 92 mm wide.
\begin{tikzpicture}[x=1mm, y=1mm, font=\scriptsize\sffamily,
    >={Stealth[length=1.6mm]},
    bb/.style={draw, fill=ln-box, align=left, inner sep=2pt, font=\scriptsize},
    nm/.style={font=\scriptsize\sffamily\bfseries, inner sep=1pt}]
  \node[text=ln-muted] (en) at (45,66) {ENTRY};
  \node[bb] (b1) at (45,57) {$d_1$: \texttt{i = 0}\\$d_2$: \texttt{s = 0}};
  \node[bb, fill=ln-accent!15] (b2) at (45,42) {$d_3$: \texttt{t = x[i]}};
  \node[bb] (b3) at (22,27) {$d_4$: \texttt{s = s + t}};
  \node[bb] (b4) at (68,27) {$d_5$: \texttt{t = -t}\\$d_6$: \texttt{s = s + t}};
  \node[bb] (b5) at (45,12) {$d_7$: \texttt{i = i + 1}};
  \node[text=ln-muted] (ex) at (45,1) {EXIT};
  \node[nm, anchor=east] at (b1.west) {B1\,};
  \node[nm, anchor=east] at (b2.west) {B2\,};
  \node[nm, anchor=east] at (b3.west) {B3\,};
  \node[nm, anchor=west] at (b4.east) {\,B4};
  \node[nm, anchor=east] at (b5.west) {B5\,};
  \draw[->] (en) -- (b1);
  \draw[->] (b1) -- (b2);
  \draw[->] (b2) -- node[above left, inner sep=1pt] {\texttt{t > 0}} (b3);
  \draw[->] (b2) -- node[above right, inner sep=1pt] {\texttt{t <= 0}} (b4);
  \draw[->] (b3) -- (b5);
  \draw[->] (b4) -- (b5);
  \draw[->] (b5) -- node[right, inner sep=1.5pt] {\texttt{i >= n}} (ex);
  \draw[->, ln-caution, thick] (b5.east) -- ++(40,0) |- node[pos=0.25, right, inner sep=1.5pt]
    {\texttt{i < n}} (b2.east);
\end{tikzpicture}

% Width: body (120 mm); the drawing is about 92 mm wide.
\begin{tikzpicture}[x=1mm, y=1mm, font=\scriptsize\sffamily,
    >={Stealth[length=1.6mm]},
    bb/.style={draw, fill=ln-box, align=left, inner sep=2pt, font=\scriptsize},
    nm/.style={font=\scriptsize\sffamily\bfseries, inner sep=1pt}]
  \node[text=ln-muted] (en) at (45,66) {ENTRY};
  \node[bb] (b1) at (45,57) {$d_1$: \texttt{i = 0}\\$d_2$: \texttt{s = 0}};
  \node[bb, fill=ln-accent!15] (b2) at (45,42) {$d_3$: \texttt{t = x[i]}};
  \node[bb] (b3) at (22,27) {$d_4$: \texttt{s = s + t}};
  \node[bb] (b4) at (68,27) {$d_5$: \texttt{t = -t}\\$d_6$: \texttt{s = s + t}};
  \node[bb] (b5) at (45,12) {$d_7$: \texttt{i = i + 1}};
  \node[text=ln-muted] (ex) at (45,1) {EXIT};
  \node[nm, anchor=east] at (b1.west) {B1\,};
  \node[nm, anchor=east] at (b2.west) {B2\,};
  \node[nm, anchor=east] at (b3.west) {B3\,};
  \node[nm, anchor=west] at (b4.east) {\,B4};
  \node[nm, anchor=east] at (b5.west) {B5\,};
  \draw[->] (en) -- (b1);
  \draw[->] (b1) -- (b2);
  \draw[->] (b2) -- node[above left, inner sep=1pt] {\texttt{t > 0}} (b3);
  \draw[->] (b2) -- node[above right, inner sep=1pt] {\texttt{t <= 0}} (b4);
  \draw[->] (b3) -- (b5);
  \draw[->] (b4) -- (b5);
  \draw[->] (b5) -- node[right, inner sep=1.5pt] {\texttt{i >= n}} (ex);
  \draw[->, ln-caution, thick] (b5.east) -- ++(40,0) |- node[pos=0.25, right, inner sep=1.5pt]
    {\texttt{i < n}} (b2.east);
\end{tikzpicture}
